\centerline{\bf \Huge ДЗ$\cdot$Матметоды}
\title{ДЗ$\cdot$Матметоды}

\begin{flushright}
    \scriptsize $-$Шурик, вы комсомолец?\\$-$Да.\\$-$Это же не наш метод.\\
    Хулиган Федя и Шурик. Операция «Ы» и другие приключения Шурика 
\end{flushright}

\noindent\llap{1.\space}В Республике Калмыкия зарегистрировано 2023 животноводческих стоянок. Докажите, что найдется такой район республики, в котором зарегистрировано более 155 животноводческих стоянок.

\noindent\llap{2.\space}Однажды Очир нарисовал гуашью прекрасную картину на небольшом прямоугольном листке бумаги. Конечно же, ему захотелось повесить ее на стену, и поэтому он поместил ее в прямоугольную рамку толщиной два сантиметра. Найдите всевозможные варианты длины и ширины прямоугольной рамки, если площадь рамки оказалась в два раза больше площади картины?

\noindent\llap{3.\space}Начальник ЭлМШ владеет операцией «Ы», которая раскладывает число на сумму пяти натуральных чисел и получает новое число с помощью перемножения этих пяти чисел. Приведите алгоритм, по которому начальник за некоторое количество операций «Ы» может из числа 30 получить число 2024. 